%% Section 4 opening and 4.1: the whole-compiler theorem and the passes
%% proved by an equation (RQ3). Sources: agda/FPC/Theorems/CorrectWrtGraphModel.agda
%% (the statements), agda/FPC/proof/Correctness/GraphModel/*Correct.agda (the
%% proofs and their header comments), loc.md (Current counts; Simplify).
%% Numbers come from loc.md (loc.py --mu) and will change; update them
%% before submission.

\section{Denotational Compiler Correctness}
\label{sec:compiler-proof}

This section proves the following two theorems about the \fpc{}
compiler.

\begin{theorem}[Compiler correctness, graph model]\label{thm:compiler}
Let $M$ be a closed term of \fpc\ of type $\kw{Int}$. For every constant $c$,
$c \in \den{M}$ if and only if $c \in \den{\pass{compile}(M)}$.
Moreover, $\den{M} \neq \emptyset$ if and only if
$\den{\pass{compile}(M)} \neq \emptyset$.
\end{theorem}

\begin{theorem}[Compiler correctness, Scott domains]\label{thm:scott}
Assume propositional truncation, function extensionality, and
propositional extensionality. Let $M$ be a closed term of \fpc\ of
type $\kw{Int}$. For every constant $c$,
$\den{M}\rho_0 = \up{(c, ())}$ if and only if
$\den{\pass{compile}(M)} = \up{(c, ())}$. Moreover, $\den{M}\rho_0$ is
defined if and only if $\den{\pass{compile}(M)}$ is defined.
\end{theorem}

Theorems~\ref{thm:compiler} and~\ref{thm:scott}, whose proofs we
mechanize in Agda, answer RQ1 affirmatively. The same statements hold
for closed terms of every type; we state them for $\kw{Int}$, the type
of a whole program's result, to align them with the operational
correctness results (Section~\ref{sec:evaluation}). Each proof is a
composition of one correctness lemma per pass.  We first show how most
of the passes can be proved correct by equational reasoning
(Section~\ref{sec:equational}). We then prove that the \pass{split
  closure} pass is correct using a logical relation
(Section~\ref{sec:approx-lr}). Finally, we sketch the corresponding
proof of correctness for Scott domains
(Section~\ref{sec:scott-delay}).

\subsection{Passes Proved by Equational Reasoning}
\label{sec:equational}

For seven of the eight passes, correctness is an equation between the
meaning of a program and the meaning of its translation.

\begin{lemma}[Passes proved by an equation]\label{lem:equational}
Let $p$ be a pass other than \pass{split closure}, and let $M$ be a
closed term of the source language of $p$. Then
$\den{M} = \den{p(M)}$.
\end{lemma}

Each proof is by induction on the term, generalizing
Lemma~\ref{lem:equational} to open terms in a source and a target
environment in which each variable means the same as its translation,
with side conditions such as nonemptiness and consistency. Every semantics is
compositional and the passes translate most forms homomorphically, so
most cases apply the same operator to equal arguments. The real work
is in the few forms that a pass changes, where an equation about the
operators of the graph model turns the source form into the target
form. Below we give these calculations for \pass{simplify} and for the
last step of \pass{enclose}. Appendix~\ref{app:calculations} gives the rest, including
\pass{select instructions}, whose key case turns a call into
$\kw{fst}$, $\kw{snd}$, $\kw{movq}$, and $\kw{callq}$.
Where a step needs a nonempty set, we show the nonempty case;
otherwise both sides are empty.

\paragraph{Simplify}

For each variable, \pass{simplify} records the term that replaces it
and, when it is bound to a statically known function, that function as
a $\lambda$-abstraction $\langle \lambda x.\, N, E_f \rangle$ paired
with the environment $E_f$ where it was defined
(Appendix~\ref{app:pass-simplify}). The main lemma for
\pass{simplify} carries the invariant that these records are true of
the source environment $\rho$ and the target environment $\rho'$. If
\pass{simplify} replaces a variable by the term $M'$, the variable
means $\den{M'}\rho'$, and if \pass{simplify} records it as the
function $\langle \lambda x.\, N, E_f \rangle$, it means
$\den{\lambda x.\, N}\rho_f$ for some environment $\rho_f$ of which
$E_f$ is true. The key case is inlining. Suppose $L$ is statically known to be the function
$\langle \lambda x.\, N, E_f \rangle$, the operand $M$ simplifies to a
small value, and the inlined body simplifies to $N'$. Then
{\small
\[
\begin{array}{@{}r@{\;}c@{\;}l@{\quad}l@{}}
  \den{L\;M}\rho & = & \den{L}\rho \app \den{M}\rho & \text{(meaning of application)}\\
    & = & \den{\lambda x.\,N}\rho_f \app \den{M}\rho & \text{(invariant for $L$)}\\
    & = & \den{N}\rho_f[x{:=}\den{M}\rho] & \text{(Lemma~\ref{lem:beta}, part 1)}\\
    & = & \den{N'}\rho' & \text{(induction hypothesis for $N$)}
\end{array}
\]}
The use of Lemma~\ref{lem:beta} needs $\den{M}\rho$ to be nonempty,
which holds because a syntactic value denotes a nonempty set, and
consistent, which holds because $\rho$ is. A $\beta$-redex
$(\lambda x.\, N)\;M$ is the same calculation without the second step.


\paragraph{Enclose}

The \pass{enclose} pass works in three steps, $\mathit{capture}$,
$\mathit{prune}$, and $\mathit{concretize}$ (\S\ref{app:pass-enclose}).
The first two produce closures of an intermediate language, Clos,
written $\kw{clos}\;(\vec{x}, y.\,N)\;\vec{M}$, which bind a name
$x_i$ for each captured value $M_i$. The last step, $\mathit{concretize}_\varphi$,
stores the captured values in a tuple. Each variable $x$ satisfies $\den{x}\rho = \den{\varphi(x)}\rho'$, where
its reference $\varphi(x)$ is either $x$ or $\kw{get}_i\,z$. Let $D_i = \den{M_i}\rho$ and
$D'_i = \den{\mathit{concretize}_\varphi(M_i)}\rho'$, so that $D_i = D'_i$,
and let $N' = \mathit{concretize}_{\varphi'}(N)$:
{\small
\[
\begin{array}{@{}r@{\;}c@{\;}l@{\quad}l@{}}
  \den{\kw{clos}\;(\vec{x}, y.\,N)\;\vec{M}}\rho
    & = & \mathsf{guard}(\vec{D};\; \Lambda(\mlam{Y} \den{N}\rho_0[\vec{x}{:=}\vec{D}, y{:=}Y])) & \text{(meaning of a Clos closure)}\\
    & = & \Lambda(\mlam{Y} \den{N}\rho_0[\vec{x}{:=}\vec{D}, y{:=}Y]) & \text{(each $D_i$ is nonempty)}\\
    & = & \Lambda(\mlam{Y} \den{N'}\rho_0[z{:=}\mathcal{T}(\vec{D}'), y{:=}Y]) & \text{(induction hypothesis for $N$)}\\
    & = & \Lambda(\mlam{X} \Lambda(\mlam{Y} \den{N'}\rho_0[z{:=}X, y{:=}Y])) \app \mathcal{T}(\vec{D}') & \text{(Lemma~\ref{lem:beta}, part 1)}\\
    & = & \den{\kw{clos}\;(z, y.\,N')\;\vec{M}'}\rho' & \text{(meaning of an EnvClos closure)}
\end{array}
\]}
The induction hypothesis applies because the code reads $x_i$ as
$\kw{get}_i\,z$, and $\pi_{\mathsf{tup}_i}(\mathcal{T}(\vec{D}')) =
D'_i = D_i$ when every $D'_i$ is nonempty.
Lemma~\ref{lem:beta} needs the tuple $\mathcal{T}(\vec{D}')$ to be
nonempty and consistent.

Each interesting case of the seven passes is a short chain of
equations like these, which is why their denotational proofs are small
(RQ3, Section~\ref{sec:evaluation}). The exception is \pass{split
  closure}, whose correctness is not an equation
(Section~\ref{sec:approx-lr}).

%% Sections 4.2 and 4.3: split closure (delay), the approximant logical relation, and the Scott-domain proof. First draft.
%% Budget: 2.75 pages (notes/paper-outline.md). The definitions follow the
%% Agda: agda/FPC/Compiler/Split.agda (the pass),
%% agda/FPC/Semantics/GraphModel/EnvClos.agda and PairClos.agda (the semantics),
%% agda/FPC/proof/Correctness/GraphModel/SplitRel.agda (the relation),
%% SplitReflect.agda (the fundamental lemma, collect),
%% SplitPreserve.agda (the forward relation) and SplitCorrect.agda (the
%% theorem). The history of the relation is told as a tour of relation
%% shapes, without names or dates (notes/history-outline.md).
%% Paper names of the languages (notes/paper-outline.md, section 2):
%% The Agda module names match the paper's language names.


\subsection{Correctness of \pass{split closure} and the Approximant Logical Relation}
\label{sec:approx-lr}

The \pass{split closure} pass separates the code of each closure from its
environment, making the closure a pair. Its proof is the hardest part of the compiler proof. This section
shows why semantic reflection for \pass{split closure} resisted our earlier
attempts and how the approximant logical relation proves it.

In EnvClos, the source language of \pass{split closure}, a closure
$\kw{clos}\;(z, y.\,N)\;\vec{M}$ still means a function. Its code $N$ is closed,
with free variables $z$ for the environment and $y$ for the argument,
and its meaning is a curried function of the environment $X$ and then
the argument $Y$. The meaning of the closure is this function applied
to the tuple of the captured values $\vec{M}$. An application $L\;M$
applies a function:
\[
\begin{array}{rcl}
    \den{\kw{clos}\;(z, y.\,N)\;\vec{M}}\rho & = &
    \Lambda(\mlam{X} \Lambda(\mlam{Y} \den{N}\rho_0[z{:=}X, y{:=}Y])) \;\app\; \mathcal{T}(\den{\vec{M}}\rho)\\
  \den{L\;M}\rho & = & \den{L}\rho \app \den{M}\rho
\end{array}
\]

In PairClos, the target language, a closure is
a pair of code and environment, and a call takes the pair apart. The
meaning of $\kw{code}\;(z, y.\,N)$ is the same curried function as before. The
meaning of $\kw{pair}\;L\;M$ is the set of pairs $\pairv{u}{v}$ with
$u$ from the code $L$ and $v$ from the environment $M$. A call applies
the first components of the pairs to their second components and then
to the argument:
\[
\begin{array}{rcl}
    \den{\kw{code}\;(z, y.\,N)}\rho & = & \Lambda(\mlam{X} \Lambda(\mlam{Y} \den{N}\rho_0[z{:=}X, y{:=}Y]))\\
  \den{\kw{pair}\;L\;M}\rho & = &
    \{\, \pairv{u}{v} \mid u \in \den{L}\rho,\; v \in \den{M}\rho \,\}\\
  \den{\kw{callc}\;L\;M}\rho & = & (\car(\den{L}\rho) \app \cdr(\den{L}\rho)) \app \den{M}\rho
\end{array}
\]
The \pass{split closure} pass translates each closure into a pair and each
application into a call. It translates every other construct
homomorphically:
\[
\begin{array}{rcl}
    \pass{split}(\kw{clos}\;(z, y.\,N)\;\vec{M}) & = &
    \kw{pair}\;(\kw{code}\;(z, y.\,\pass{split}(N)))\;(\kw{tuple}\;\pass{split}(\vec{M}))\\
  \pass{split}(L\;M) & = &
    \kw{callc}\;\pass{split}(L)\;\pass{split}(M)
\end{array}
\]

\begin{figure}[t]
\small\figdefs
\defhead{}{v \sqsubseteq D'}
\vspace{-10pt}
\[
\begin{array}{@{}r@{\;}c@{\;}l@{}}
  (U \mapsto w) \sqsubseteq D' & = &
    (\exists u, t.\; \pairv{u}{t} \in D') \;\land\;
    \forall E'.\; \mathsf{con}(E') \land U \sqsubseteq^* E'
      \Rightarrow w \sqsubseteq (\car(D') \app \cdr(D')) \app E'\\
  \nu \sqsubseteq D' & = & \exists u, t.\; \pairv{u}{t} \in D'\\
  t(v) \sqsubseteq D' & = & v \sqsubseteq \pi_t(D')
    \qquad \text{for } t \in \{\mathsf{tup}_i, \mathsf{left}, \mathsf{right}\}\\
  v \sqsubseteq D' & = & v \in D'
    \qquad \text{for the other elements}
\end{array}
\]
\vspace{-5pt}
\defhead{}{U \sqsubseteq^* E'}
\vspace{-10pt}
\[
\begin{array}{@{}r@{\;}c@{\;}l@{}}
  [\,] \sqsubseteq^* E' & = & \mathsf{true}\\
  (u :: U) \sqsubseteq^* E' & = & u \sqsubseteq E' \land U \sqsubseteq^* E'
\end{array}
\]
\caption{The relation $v \sqsubseteq D'$ (``$v$ is included in the
  behavior of $D'$'') for the preservation half of \pass{split closure}. It
  relates a source element $v$ (of EnvClos) to a set $D'$ of target
  elements (of PairClos) and is defined by recursion on $v$. The
  relation $U \sqsubseteq^* E'$ extends it to a finite list $U$.}
\label{fig:preserve}
\end{figure}

\paragraph{Why preservation is easy and reflection is not}

The correctness of \pass{split closure} is not an equation, because
the pass changes the denotation of a function. In EnvClos, a closure
denotes the graph of the function. In PairClos, it denotes pairs of
the graph of the function curried over the values of its free
variables and the tuple of those values. Lemma~\ref{lem:delay}
therefore states only that a closed program and its translation
produce the same constants and that one produces something exactly
when the other does. The forward half, preservation, requires only a
relation between a \emph{single} source element and the \emph{set} of
target elements, defined by recursion on the source element
(Figure~\ref{fig:preserve}). We read $v \sqsubseteq D'$ as ``$v$ is
included in the behavior of $D'$''. The source element is finite, so
the recursion is structural, and preservation is a routine induction
on $M$.

The backward half, reflection, must show that the source accounts for
every element that the target produces. Here the representation of
closures gets in the way. Consider the function $\lambda y.\, x + y$
in an environment where $x$ is $9$. In EnvClos, its meaning 
includes entries $[1] \mapsto 10$ and $[6] \mapsto 15$. After
\pass{split closure}, its code takes the environment as an extra
parameter. The meaning of the code is a table with entries for
\emph{every} environment, such as the following two:
\[
  [9] \mapsto ([1] \mapsto 10) \qquad
  [2] \mapsto ([1] \mapsto 3)  \qquad \ldots
\]
The program pairs this table only with environment $[9]$, so the entry
for $[2]$ does not correspond to a source function. Worse, the
denotation of a call takes the table from $\car$ and the environment
from $\cdr$ \emph{separately}, so it may combine the table of one pair
with the environment of another.

\paragraph{Design alternatives: which sides to relate}

Before we define our relation for reflection, we explain how we chose
its shape, because the obvious designs fail. A relation for reflection
relates a source side to a target side, and each side can be a single
graph element, a finite list of elements, or a set that may be
infinite. Of these nine designs, the following reasons rule out all
but one (Appendix~\ref{app:alr-example} gives
examples and the variants we tried in Agda).
\begin{itemize}
\item \emph{A single target element}, against any source side, fails
  because a call combines the code entry of one pair with the
  environment of another, so no single pair accounts for its result.
\item \emph{A single source element or a finite source list}, against
  any target side, fails in the function case, because infinitely many related
  source argument sets have no element in common, so no finite source
  side has an entry for each of them.
\item \emph{Whole denotations on both sides} handle the mixing, but the
  function case uses the relation in negative position, where nothing
  decreases. A well-founded definition would need a later modality or
  a construction by stages, as in our Scott-domain proof
(Section~\ref{sec:scott-delay}).
\item \emph{A target set of bounded depth} against the source
  denotation is well founded, but the target denotation is unbounded,
  and the outputs of a call can come from code entries of any depth.
\end{itemize}
The remaining design, which we adopt, relates a finite list of target
elements to the whole source denotation. We also had to choose what
the definition recurses on. Recursion on types, as in the proofs of
\citet{Chlipala:2007fk} and \citet{Nielsen:2000aa}, is not an option.
Because $\kw{fold}$ and $\kw{unfold}$ mean the identity, the relation at
$\mu X.\,A$ would be the relation at $A[\mu X.\,A/X]$, which is not a smaller type. Our relation instead recurses on the depth of the target elements.

\paragraph{The key insight}

When a call $(\car(D') \app \cdr(D')) \app E'$ produces a result $w$, it uses only
finitely many pairs of $D'$. One supplies a code entry
$U_e \mapsto (U_a \mapsto w)$, and others supply the environments in
$U_e$. These pairs form a finite list $V \subseteq D'$ that contains
both halves of the call, so the relation matches the code and the environments \emph{inside} $V$. The relation lets a call use exactly the code
entries whose environments $V$ itself provides, which excludes the
junk. The finite list also makes the relation well founded.

\begin{figure}[t]
\small\figdefs
\defhead{Usable entries}{\entryin{U_a}{w}{V}}
\[
\begin{array}{@{}r@{\;}c@{\;}l@{}}
  \entryin{U_a}{w}{V} & \iff &
    \exists\, U_e.\;
    (U_e \mapsto (U_a \mapsto w)) \in \car(V)
    \;\land\; U_e \neq [\,]
    \;\land\; U_e \subseteq \cdr(V)
\end{array}
\]

\defhead{Arguments of usable entries}{\mathsf{args}(V) = D}
\[
  \mathsf{args}(V) \;=\; \{\, u \mid \entryin{U_a}{w}{V},\; u \in U_a \,\}
\]

\defhead{Self-contained calls}{\callsto{V}{U}{W}}
\[
  \callsto{V}{U}{W} \iff
    U \subseteq \mathsf{args}(V) \;\land\;
    (\forall w \in W.\; \exists U_a.\; \entryin{U_a}{w}{V}
      \land U_a \neq [\,] \land U_a \subseteq U)
\]

\defhead{The approximant logical relation}{D \ALR_k V \qquad D \ALR V}
\[
\begin{array}{@{}r@{\;}c@{\;}l@{}}
  D \ALR_0 V & \iff & \mathsf{true}\\
  D \ALR_{k+1} V & \iff &
    (\forall c.\; c \in V \Rightarrow c \in D)\\
  & \land & (V \neq [\,] \Rightarrow D \neq \emptyset)\\
  & \land & (\forall U, W, E.\;
     \callsto{V}{U}{W} \land \mathsf{con}(E) \land \mathsf{con}(U)
     \land E \ALR_k U \Rightarrow (D \app E) \ALR_k W)\\
  & \land & (\forall \pi, W.\; W \subseteq \pi(V) \Rightarrow \pi(D) \ALR_k W)\\[4pt]
  D \ALR V & \iff & D \ALR_{1 + \depth(V)} V
\end{array}
\]
\caption{The approximant logical relation $D \ALR V$ between a set $D$
  of source (EnvClos) elements and a finite list $V$ of target (PairClos)
  elements. Here $c$
  ranges over constants and $\pi$ over the projections of tuple
    components and sum injections; $\car$ and $\cdr$ are as in
  Figure~\ref{fig:graph-ops}, $\mathsf{con}$ is consistency
  (Section~\ref{sec:graph-models}), and $\depth(V)$ is the largest
  depth of an element of $V$ (Appendix~\ref{app:depth}).}
\label{fig:alr}
\end{figure}

\paragraph{The approximant logical relation}

Figure~\ref{fig:alr} defines the approximant logical relation
$D \ALR V$, which relates a finite list $V$ of target elements to a
source set $D$. We read $D \ALR V$ as ``$V$ approximates $D$''. Its building block is the \emph{usable entry} of a
list. A target closure is a set of pairs $\pairv{u}{t}$ of code and
environment. The code $u$ is a curried function, so its entries have
the form $U_e \mapsto (U_a \mapsto w)$, which maps the environment $U_e$
to a function that maps the arguments $U_a$ to $w$. We write
$\entryin{U_a}{w}{V}$, read ``$U_a \mapsto w$ is a usable entry of
$V$'', when such an entry is among the codes $\car(V)$
\emph{and} the environments $U_e$ are among the environments
$\cdr(V)$ that $V$ itself provides. The second condition makes the
relation sound in the presence of junk. A code entry for an environment
that the program never supplies, such as the entry for $[2]$ in the
example above, is not usable. A call on $D'$ takes the code from $\car(D')$ and the environment
from $\cdr(D')$. In the same way,
the code and the environment of a usable entry may come from different
pairs of $V$.

We write $\callsto{V}{U}{W}$, read ``$V$ maps $U$ to $W$'', when
every argument in $U$ is taken by some usable entry of $V$ (the first
conjunct, $U \subseteq \mathsf{args}(V)$) and each result $w$ in $W$
is produced by a usable entry of $V$ whose arguments are among $U$. We
call this a \emph{self-contained call}, because $V$ supplies both the
code entries and the environments.

The relation $D \ALR_{k+1} V$ has four clauses, which use ${\ALR_k}$
at the recursive occurrences:
\begin{enumerate}
  \item every constant in $V$ is in $D$;
  \item if $V$ is nonempty, then $D$ is nonempty;
  \item whenever $V$ maps $U$ to $W$ and $E$ is a consistent set
    related to $U$, the set $D \app E$ is related to $W$; and
  \item every sublist of a projection of $V$ is related to the same
    projection of $D$.
\end{enumerate}
Appendix~\ref{app:alr-example} works clause~3 through an example.



\paragraph{Termination of the Approximant Logical Relation}

The relation $D \ALR V$ is defined by recursion on the depth of
$V$. Every recursive occurrence, including the negative one at the
arguments of a call, is at a list shallower than $V$. The arguments $U$ are shallower than $V$
because the first conjunct of a self-contained call requires $U \subseteq
\mathsf{args}(V)$; without that conjunct, $U$ could contain elements
deeper than $V$. The results $W$ are shallower than $V$ because each
of them is the output of an entry in $V$. Finally, the projections of
$V$ are shallower than $V$ because they consist of components of the
elements of $V$.

The index $k$ of ${\ALR_k}$ serves only to ensure termination. Any
index that bounds the depth of $V$ gives the same relation, and
therefore $D \ALR V$ is equivalent to the four clauses with ${\ALR}$
itself at the recursive occurrences
(Lemma~\ref{lem:depth-irrelevant}), and the proofs never mention
the index.

\begin{lemma}[Depth irrelevance and unfolding]\label{lem:depth-irrelevant}
\leavevmode
\begin{enumerate}
\item If every element in $V$ has depth less than $k$ and less than
  $j$, then $D \ALR_k V$ if and only if $D \ALR_j V$.
\item $D \ALR V$ if and only if the four clauses of
  $D \ALR_{k+1} V$ hold with ${\ALR}$ in place of ${\ALR_k}$.
\end{enumerate}
\end{lemma}

The relation is also monotone in $D$ and antitone in $V$, that is,
enlarging the source set or taking a sublist of the target list preserves it.

\paragraph{The fundamental lemma}

A source environment $\rho$ and a target environment $\rho'$ are
related, written $\rho \ALR \rho'$ and read ``$\rho'$ approximates
$\rho$'', when $\rho(x) \ALR V$ for every
variable $x$ and every $V \subseteq \rho'(x)$. We also require that both environments be
consistent and that $\rho$ map every variable to a nonempty set.

\begin{lemma}[Fundamental lemma]\label{lem:delay-fundamental}
Let $M$ be a term of EnvClos. If $\rho \ALR \rho'$ and
$V \subseteq \den{\pass{split}(M)}\rho'$, then
$\den{M}\rho \ALR V$.
\end{lemma}

The proof is by induction on $M$. Throughout, the second part of
Lemma~\ref{lem:depth-irrelevant} lets us pass between $D \ALR V$ and
its four clauses, with ${\ALR}$ at the recursive occurrences. The case of a variable follows from $\rho \ALR \rho'$, and the cases of $\kw{fold}$, $\kw{unfold}$, tuples, injections, $\kw{let}$, and $\kw{case}$ follow from the induction hypotheses for their subterms. Two cases carry the weight.

\fbox{\emph{Case $M = L\;K$.}} Then
$\pass{split}(M) = \kw{callc}\,\pass{split}(L)\,\pass{split}(K)$,
so $V \subseteq (\car(D') \app \cdr(D')) \app E'$ for
$D' = \den{\pass{split}(L)}\rho'$ and
$E' = \den{\pass{split}(K)}\rho'$. Each element of $V$ is produced by
a code entry from $\car(D')$, finitely many environments from
$\cdr(D')$, and finitely many arguments from $E'$. We collect the pairs
of $D'$ that supply these codes and environments into a finite list
$W \subseteq D'$, and the arguments into a finite list
$U \subseteq E'$. Then $W$ maps $U$ to $V$. The induction hypotheses
give $\den{L}\rho \ALR W$ and $\den{K}\rho \ALR U$, and both
$\den{K}\rho$ and $U$ are consistent. Clause~3 for
$\den{L}\rho \ALR W$, applied to this self-contained call with
$E = \den{K}\rho$, gives $\den{L}\rho \app \den{K}\rho \ALR V$. That is,
$\den{M}\rho \ALR V$.

\fbox{\emph{Case $M = \kw{clos}\,(z, y.\,N)\,\vec{K}$.}} Then
$\pass{split}(M) = \kw{pair}\,(\kw{code}\,(z, y.\,\pass{split}(N)))\,(\kw{tuple}\,\pass{split}(\vec{K}))$,
so $V$ is a list of pairs whose environments lie in
$T' = \mathcal{T}(\den{\pass{split}(\vec{K})}\rho')$. Let
$D = \den{M}\rho$ and $T = \mathcal{T}(\den{\vec{K}}\rho)$. We prove the
four clauses of $D \ALR V$. Clauses 1 and 4 hold trivially
because $V$ contains only pairs.

For clause 2, suppose $\pairv{u}{t'} \in V$. Then
$t' \in T'$, so $v_i \in \den{\pass{split}(K_i)}\rho'$ for some $v_i$ at each $i$.
The induction hypothesis for $K_i$, with $\rho \ALR \rho'$ and
$[v_i] \subseteq \den{\pass{split}(K_i)}\rho'$, gives
$\den{K_i}\rho \ALR [v_i]$, and clause 2 of the $\ALR$ relation makes
$\den{K_i}\rho$ nonempty. Hence $t \in T$ for some $t$. Every $\Lambda(F)$
contains $\nu$, so the code of the source closure contains the entry
$[t] \mapsto \nu$, and applying that to $t \in T$ gives $\nu \in D$.

Clause 3 is the heart of the proof. Suppose that $V$ maps $U$ to $W$
and $E \ALR U$. The bodies of the source and target closures run in
the environments $\rho_N = \rho_0[z{:=}T, y{:=}E]$ and
$\rho'_N = \rho_0[z{:=}\cdr(V), y{:=}U]$. We show the two premises of
the induction hypothesis for the body $N$.
First, $W \subseteq \den{\pass{split}(N)}\rho'_N$. Each $w \in W$ comes
from a code entry $U_e \mapsto (U_a \mapsto w)$ in $\car(V)$ with
$U_e \subseteq \cdr(V)$ and $U_a \subseteq U$, so
$w \in \den{\pass{split}(N)}\rho_0[z{:=}U_e, y{:=}U_a]
\subseteq \den{\pass{split}(N)}\rho'_N$ by monotonicity
(Lemma~\ref{lem:monotone}).
Second, $\rho_N \ALR \rho'_N$. We have $T \ALR \cdr(V)$ by the induction
hypothesis for $\vec{K}$, because $\cdr(V) \subseteq T'$, and
$E \ALR U$.
With the two premises in hand, the induction hypothesis for $N$
gives us $\den{N}\rho_N \ALR W$.
Finally, Lemma~\ref{lem:beta} yields $\den{N}\rho_N \subseteq D \app E$,
and ${\ALR}$ is monotone in its source set, so $D \app E \ALR W$.

\begin{lemma}[Correctness of \pass{split closure}]\label{lem:delay}
Let $M$ be a closed term of EnvClos. For every constant $c$,
$c \in \den{M}$ if and only if $c \in \den{\pass{split}(M)}$.
Moreover, $\den{M} \neq \emptyset$ if and only if
$\den{\pass{split}(M)} \neq \emptyset$.
\end{lemma}

Because $\rho_0 \ALR \rho_0$, reflection follows from
Lemma~\ref{lem:delay-fundamental}. If $c \in \den{\pass{split}(M)}$,
the lemma applied to the one-element list $[c]$ gives
$\den{M} \ALR [c]$, and clause~1 gives $c \in \den{M}$. If
$\den{\pass{split}(M)}$ contains an element $v$, the lemma applied
to $[v]$ and clause~2 show that $\den{M}$ is nonempty. Preservation
follows from the relation $\sqsubseteq$ of
Figure~\ref{fig:preserve}. For a constant, $c \sqsubseteq D'$ means
$c \in D'$, and $v \sqsubseteq D'$ implies $D' \neq \emptyset$.

\subsection{Proof of Compiler Correctness with Scott Domains}
\label{sec:scott-delay}

We also proved the compiler correct with Scott domains
(Appendix~\ref{app:scott}). There a value is an element of a
universal domain $\mathbb{D}$ of tagged values, divergence is its least
element $\bot$, and \pass{split closure} is again
the only pass that needs a relation. For space, this section only
summarizes that proof; Appendix~\ref{app:scott-delay} gives it in
full. We define the relation by
recursion on the stages $\mathbb{D}_0, \mathbb{D}_1, \ldots$ of the
bilimit that constructs $\mathbb{D}$, as \citet{Reynolds:1974aa} did.
Let $d$ be a source value and $d'$ a target value, and write
$\pi^\infty_n(d) \in \mathbb{D}_n$ for the approximation of $d$ at
stage $n$. For preservation, the stage relation
$\pi^\infty_n(d) \SR{+}{n} d'$ relates the approximation of $d$ to the
whole of $d'$. We write $d \SRel{+} d'$, read ``$d'$ preserves $d$'',
when it holds at every stage. For reflection, the roles swap, and
$d' \SRel{-} d$, read ``$d$ reflects $d'$'', holds when
$\pi^\infty_n(d') \SR{-}{n} d$ at every stage
(Figure~\ref{fig:scott-delay} in Appendix~\ref{app:scott-delay}).
The clauses for closures are
\[
\begin{array}{@{}r@{\;}c@{\;}l@{}}
  (\mathit{fun}, h) \mathrel{\mathcal{M}^{+}_n} (\mathit{pr}, (c, e))
    & \iff & \forall v, a.\;
      v \SR{+}{n} a \Rightarrow h(v) \SR{+}{n} \mathit{call}(c, e, a)\\
  (\mathit{pr}, (c, e)) \mathrel{\mathcal{M}^{-}_{m+2}} (\mathit{fun}, g)
    & \iff & \forall v, a.\;
      v \SR{-}{m} a \Rightarrow \mathit{call}_m((c, e), v) \SR{-}{m} g'(a)
\end{array}
\]
where $\mathit{call}(c, e, a)$ applies the code $c$ to the environment
$e$ and then to the argument $a$. The preservation clause has the
shape of the relation that \citet{Nielsen:2000aa} used for
defunctionalization. The reflection clause relates calls across
stages, because a call of a pair at stage $m + 2$ applies its code
twice and shows its result only at stage $m$
($\mathit{call}_m$). Because the closure clause looks three levels
into the target domain but only one into the source, the method of
\citet{Pitts:1996aa} does not apply directly. The result,
Lemma~\ref{lem:scott-delay}, is the counterpart of
Lemma~\ref{lem:delay} and the key step in proving
Theorem~\ref{thm:scott}.


